\hypertarget{appendix-313ee7a0ba7b}{}
\section*{Memory-mapped files}
\section{Memory-mapped files}\label{memory-mapped-files}

\textbf{Curriculum depth:} Introduced · \textbf{First seen:} C003 · \textbf{Latest development:} C003

\subsection{Definition and mechanism}\label{definition-and-mechanism}

A file mapping makes file-backed bytes accessible through a virtual region. Fault handling can supply pages lazily; ownership, access flags, mapping length, and synchronization still matter. A mapped byte buffer is not automatically a set of live aligned C++ objects. File persistence and virtual access are separate mechanisms.

\subsection{Prerequisites}\label{prerequisites}

\begin{itemize}
\tightlist
\item
  \hyperref[appendix-cc0909a7ef14]{Page faults and demand paging} --- Deferred private copying is implemented through a recoverable write fault.
\item
  \href{}{Ownership and non-owning views (05-Engine-I, appendix-d52e4486f527)} --- A mapping or publication pipeline must keep referenced bytes valid.
\end{itemize}

\subsection{Used by}\label{used-by}

Update downstream uses when later chunks build on this concept.

\subsection{Recall and next depth}\label{recall-and-next-depth}

Separate mapping setup, first-touch loading, typed parsing and persistence.

\textbf{Pending:} Platform mapping APIs, sharing modes, truncation and synchronization.

\subsection{Source and navigation}\label{source-and-navigation}

\href{Continuous\%20Notes/Computing\%20Foundations\%20-\%20From\%20Bits\%20to\%20C++\%20and\%20Systems\%20APIs\#Where\%20the\%20bytes\%20live}{C003 teaching chapter} · \hyperref[appendix-fd090d982a15]{Part definitions} · \href{}{Dependency Map (02-Foundations-II, appendix-4ad47bdd269a)}
